% PAGES: 57-59

\subsection{Finding discount factors using forward substitution}

The value of a bond is equal to the sum of the present values of all its future cash
flows. Therefore, known prices of several bonds provide information regarding the
discount factors corresponding to all the cash flow dates of the bonds.

If the number of cash flow dates is higher than the number of bonds, then there is not
enough information to uniquely determine the discount factors. If the number of bonds
is higher than the number of cash flow dates, there is a redundancy of data and
arbitrage opportunities may exist.

However, if the number of cash flow dates is equal to the number of bonds, and if
there are no redundancies (i.e., if no bond can be synthesized by taking positions in
the other bonds), then it is possible to uniquely identify the discount factors
corresponding to the cash flow dates (and the corresponding values of the zero rates)
from the bond prices.

In this section, we present an example where the discount factors can be computed
using forward substitution. A more general example requiring LU decomposition for
computing the discount factors can be found in section 2.7.2.

Recall that the value of a bond is equal to the sum of the present value of all its
future cash flows. In other words, if $B$ is the value of a bond with future cash
flows $c_i$ to be paid at times $t_i$, and if $\text{Disc}(t_i)$ are the discount
factors corresponding to time $t_i$, for $i=1:n$, then
\begin{equation}
B \;=\; \sum_{i=1}^n c_i\text{Disc}(t_i).
\end{equation}

Note that the discount factor $\text{Disc}(t_i)$ at time $t_i$ is uniquely determined
by the risk-free zero rate $r(0,t_i)$ at time $t_i$. For example, if interest is
continuously compounded (which is assumed to be the case throughout the book, unless
otherwise specified), then
\begin{equation}
\text{Disc}(t_i) \;=\; e^{-t_ir(0,t_i)} \;=\; \exp\left(-t_ir(0,t_i)\right).
\end{equation}

For the example below, recall that an annual coupon bond with face value \$100,
coupon rate $C$, and maturity $T$ (measured in years) pays the holder of the bond a
coupon payment equal to $C\cdot 100$ every year, except at maturity. The final payment
at maturity $T$ is equal to the face value of the bond plus one coupon payment, i.e.,
$(1+C)100$. For example, an annual coupon bond with 15 months to maturity and coupon
rate 5\%, i.e., $C=0.05$, has two cash flow dates, in 3 and in 15 months,
corresponding to $t_1 = \frac{3}{12} = \frac{1}{4}$, and $t_2 = \frac{15}{12} =
\frac{5}{4}$, respectively. The corresponding cash flows are $c_1=5$ and $c_2=105$.

\begin{examplex}
The prices of the following annual coupon bonds with face value \$100 are given:

\begin{center}
\begin{tabular}{ccc}
Bond Maturity & Coupon Rate & Bond Price \\[0.4em]
1 year & 0 & \$98 \\
2 years & 6\% & \$104 \\
3 years & 8\% & \$111 \\
4 years & 5\% & \$102 \\
\end{tabular}
\end{center}

The cash flows and the cash flow dates of the bonds are recorded in the table below:

\begin{center}
\begin{tabular}{|l|l|}
\hline
Maturity & Cash Flow \& Date \\
\hline
1 year & \$100 in 1 year \\
\hline
2 years & \begin{tabular}[t]{@{}l@{}}\$6 in 1 year\\ \$106 in 2 years\end{tabular} \\
\hline
3 years & \begin{tabular}[t]{@{}l@{}}\$8 in 1 year\\ \$8 in 2 years\\
\$108 in 3 years\end{tabular} \\
\hline
4 years & \begin{tabular}[t]{@{}l@{}}\$5 in 1 year\\ \$5 in 2 years\\
\$5 in 3 years\\ \$105 in 4 years\end{tabular} \\
\hline
\end{tabular}
\end{center}

Note that the four bonds have exactly four cash flow dates, i.e., 1 year, 2 years, 3
years, and 4 years. For $i=1:4$, denote by $d_i$ the discount factor corresponding to
time equal to $i$ years. Using formula (2.11) for the values of the bonds, we find
that
\begin{align*}
98 &= 100d_1; \\
104 &= 6d_1 + 106d_2; \\
111 &= 8d_1 + 8d_2 + 108d_3; \\
102 &= 5d_1 + 5d_2 + 5d_3 + 105d_4.
\end{align*}

This linear system can be written in matrix notation as $Ld = b$, where $L$ is a
lower triangular matrix, and $d$ and $b$ are column vectors given by
\[
L = \begin{pmatrix} 100 & 0 & 0 & 0 \\ 6 & 106 & 0 & 0 \\ 8 & 8 & 108 & 0 \\
5 & 5 & 5 & 105 \end{pmatrix}; \quad
d = \begin{pmatrix} d_1 \\ d_2 \\ d_3 \\ d_4 \end{pmatrix}; \quad
b = \begin{pmatrix} 98 \\ 104 \\ 111 \\ 102 \end{pmatrix}.
\]

The solution of the linear system $Ld = b$ is found using forward
substitution,\footnote{Without using the routine forward\_subst, $d_1$, $d_2$, $d_3$,
and $d_4$ can be obtained as follows:
\[
d_1 \;=\; \frac{98}{100} \;=\; 0.98;
\]
\[
d_2 \;=\; \frac{104 - 6d_1}{106} \;=\; 0.9257;
\]
\[
d_3 \;=\; \frac{111 - 8d_1 - 8d_2}{108} \;=\; 0.8866;
\]
\[
d_4 \;=\; \frac{102 - 5d_1 - 5d_2 - 5d_3}{105} \;=\; 0.8385.
\]} i.e.,
\[
d \;=\; \text{forward\_subst}(L,b) \;=\; \begin{pmatrix} 0.98 \\ 0.9257 \\ 0.8866 \\
0.8385 \end{pmatrix}.
\]

Formula (2.12) can be used to find the corresponding continuously compounded zero
rates:
\begin{align*}
d_1 = \text{Disc}(1) = e^{-r(0,1)} \quad &\Longleftrightarrow \quad r(0,1) =
-\ln(d_1) = 0.0202 = 2.02\%; \\
d_2 = \text{Disc}(2) = e^{-2r(0,2)} \quad &\Longleftrightarrow \quad r(0,2) =
-\frac{\ln(d_2)}{2} = 0.0386 = 3.86\%; \\
d_3 = \text{Disc}(3) = e^{-3r(0,3)} \quad &\Longleftrightarrow \quad r(0,3) =
-\frac{\ln(d_3)}{3} = 0.0401 = 4.01\%; \\
d_4 = \text{Disc}(4) = e^{-4r(0,4)} \quad &\Longleftrightarrow \quad r(0,4) =
-\frac{\ln(d_4)}{4} = 0.0440 = 4.40\%.
\end{align*}
\end{examplex}
